\section*{Abstract}

Data storage is a fundamental component of modern computing systems, directly influencing the 
performance and efficiency of the operating system. This case study examines and compares three 
major disk storage technologies — Hard Disk Drives (HDD), Solid State Drives (SSD), and NVMe 
drives — analyzing their architecture, working principles, performance characteristics, and cost 
considerations. The study also explores how operating systems interact with each storage type 
through device drivers, I/O scheduling, and file system management. Findings indicate that while 
HDDs remain relevant for large-capacity, low-cost storage, SSDs offer a significant improvement in 
speed and durability, and NVMe drives represent the current peak of consumer storage performance due 
to their direct PCIe interface and extremely low latency. The study concludes that the choice of 
storage technology has a direct and measurable impact on overall system performance, and that modern 
operating systems must be designed to effectively manage and optimize each type of storage medium.